\documentclass[../../../include/open-logic-section]{subfiles}

\begin{document}

\olfileid[id]{his}{bio}{god}

\olsection{Kurt G\"odel}

\olphoto{goedel-kurt}{Kurt G{\"o}del}

Kurt G{\"o}del (\textsc{ger}-dle) lahir pada 28 April 1906 di Br{\"u}nn,
Kekaisaran Austria-Hungaria (kini Brno di Republik Ceko). Karena sifatnya
yang ingin tahu dan cerdas, Kurtele kecil sering dipanggil ``Der kleine Herr
Warum'' (Tuan Kecil Mengapa) oleh keluarganya. Sejak sekolah dasar ia unggul
dalam pelajaran; hanya dalam matematika nilainya pernah berada di bawah
nilai tertinggi. G{\"o}del kerap tidak masuk sekolah karena kesehatannya
buruk dan dibebaskan dari pendidikan jasmani. Semasa kanak-kanak ia
didiagnosis menderita demam rematik. Sepanjang hidupnya, ia percaya bahwa
penyakit itu merusak jantungnya secara permanen, meskipun evaluasi medis
menyatakan sebaliknya.

G{\"o}del mulai belajar di Universitas Wina pada tahun 1924 dan menyelesaikan
studi doktoralnya pada tahun 1929. Mula-mula ia bermaksud mempelajari fisika,
tetapi minatnya segera beralih ke matematika, terutama logika, sebagian
karena pengaruh filsuf Rudolf Carnap. Disertasinya, yang ditulis di bawah
bimbingan Hans Hahn, membuktikan teorema kelengkapan bagi logika predikat
orde pertama dengan identitas \citep{Godel1929}. Hanya setahun kemudian, ia
memperoleh hasil-hasilnya yang paling terkenal---teorema ketaklengkapan
pertama dan kedua (diterbitkan dalam \citealt{Godel1931}). Selama berada di
Wina, G{\"o}del sangat aktif dalam Lingkaran Wina, kelompok filsuf
berorientasi ilmiah yang mencakup Carnap, yang karyanya sangat dipengaruhi
oleh hasil-hasil G{\"o}del.

Pada tahun 1938, G\"odel menikahi Adele Nimbursky. Orang tuanya tidak senang:
Adele bukan hanya enam tahun lebih tua dan pernah bercerai, melainkan juga
bekerja sebagai penari di klub malam. Namun, tekanan sosial tidak
memengaruhi G{\"o}del, dan mereka tetap menjalani pernikahan yang bahagia
hingga kematiannya.

Setelah Jerman Nazi mencaplok Austria pada tahun 1938, G{\"o}del dan Adele
beremigrasi ke Amerika Serikat, tempat ia menerima jabatan di Institute for
Advanced Study di Princeton, New Jersey. Meskipun bersifat tertutup dan
eksentrik, masa G{\"o}del di Princeton penuh kolaborasi dan produktif. Ia
menerbitkan esai dalam teori himpunan, filsafat, dan fisika. Yang terutama
patut dicatat, ia menjalin persahabatan yang sangat erat dengan rekannya di
IAS, Albert Einstein.

Pada tahun-tahun terakhir hidupnya, kesehatan mental G{\"o}del memburuk.
Ketika istrinya dirawat di rumah sakit pada tahun 1977, ia tidak lagi dapat
memasakkan makanan untuknya. Setelah mengalami masalah kesehatan mental
sepanjang hidupnya, ia dikuasai paranoia. Karena sangat takut diracuni,
G{\"o}del menolak makan. Ia meninggal akibat kelaparan pada 14 Januari 1978
di Princeton.

\begin{reading}
Untuk biografi lengkap mengenai kehidupan G{\"o}del, lihat
\citet{Dawson1997}. Untuk tulisan biografis lebih lanjut serta esai mengenai
sumbangan G{\"o}del terhadap logika dan filsafat, lihat \citet{Wang1990},
\citet{Baaz2011}, \citet{Takeuti2003}, dan \citet{Sigmund2007}.

Tesis PhD G{\"o}del tersedia dalam bahasa Jerman aslinya
\citep{Godel1929}. Teks asli teorema-teorema ketaklengkapan terdapat dalam
\citep{Godel1931}. Semua tulisan G\"odel yang diterbitkan maupun tidak
diterbitkan, beserta pilihan korespondensi, tersedia dalam bahasa Inggris
dalam \emph{Collected Papers} \citet{Godel1986,Godel1990}.

Untuk pembahasan terperinci mengenai teorema ketaklengkapan G{\"o}del, lihat
\citet{Smith2013}. Untuk pembahasan informal dan filosofis mengenai
teorema-teorema G{\"o}del, lihat siniar Mark Linsenmayer
\citep{Linsenmayer2014}.
\end{reading}

\end{document}
